% G3-0 Grafik statistik data. Data dari scripts/grafik_data.py (lewat
% btcpred.data.load di prediksi_btc). Hitam-putih; perbarui datanya saat
% data final Februari 2027 tersedia.
\begin{tikzpicture}
  \begin{axis}[name=a, width=0.86\textwidth, height=4.6cm,
      title={(a) Harga \textit{typical} BTC/IDR setiap awal jam},
      title style={font=\footnotesize},
      xlabel={Hari ke- sejak 10 Oktober 2026 (UTC)},
      ylabel={Harga (juta Rp)},
      label style={font=\footnotesize}, tick label style={font=\footnotesize},
      /pgf/number format/.cd, use comma, 1000 sep={.},
      xmin=0, enlarge x limits=0.01, grid=none]
    \addplot[black, thin, mark=none] table[x=hari, y=harga_juta]
      {assets/data/harga_per_jam.dat};
  \end{axis}
  \begin{axis}[at={(a.south)}, anchor=north, yshift=-2.3cm,
      width=0.86\textwidth, height=4.6cm,
      title={(b) Sebaran \textit{return} logaritmik 1 menit},
      title style={font=\footnotesize},
      xlabel={\textit{Return} (bp); bin tepi memuat nilai di luar $\pm$20 bp},
      ylabel={Persentase (\%)},
      label style={font=\footnotesize}, tick label style={font=\footnotesize},
      /pgf/number format/.cd, use comma, 1000 sep={.},
      ybar, bar width=0.8, xmin=-21, xmax=21, ymin=0]
    \addplot[draw=black, fill=black!25] table[x=tengah, y=persen]
      {assets/data/hist_return_1m.dat};
  \end{axis}
\end{tikzpicture}
